% ============================================================================
%  CHAPITRE 3 -- CANAL SIMULÉ ET CHAÎNE SÉPARÉE
% ============================================================================
\chapter{Canal simulé et chaîne séparée}
\label{chap:separee}

La référence à battre doit être implémentée honnêtement, sinon la comparaison ne vaut rien.

\section{Le canal}
\label{sec:canal_impl}

Implémenté dans \srcfile{semcom/channel.py}, il enchaîne trois opérations communes aux trois systèmes.

\begin{enumerate}
  \item \textbf{Troncature} : seuls les $c$ premiers canaux latents sont émis, les autres sont mis à zéro au
        récepteur.
  \item \textbf{Normalisation} : la puissance moyenne vaut exactement 1 par symbole complexe,
        \begin{equation}
          x = z_{1:c}\sqrt{\frac{k}{\sum_i z_i^2}}, \qquad k = \tfrac{1}{2}\,\mathrm{card}(z_{1:c}) .
          \label{eq:power}
        \end{equation}
  \item \textbf{Bruit} : AWGN de variance $\sigma^2 = 10^{-\mathrm{SNR}/10}$, ou Rayleigh — tirage de
        $|h|^2$, SNR instantané $\mathrm{SNR}+10\log_{10}|h|^2$, puis même bruit.
\end{enumerate}

\section{Codage de source}
\label{sec:codage_source}

\paragraph{Grandes images : JPEG.} Sur Kodak et en vidéo, JPEG~\cite{wallace1992} à la plus haute qualité qui
tienne dans le budget. Au point $1/12$ en QPSK 1/2, une image Kodak dispose de \num{12216}~octets : la
qualité retenue est 6, pour \num{11084}~octets.

\paragraph{Petites images : auto-encodeur 4 bits.} Sur CIFAR-10, JPEG est inutilisable, et c'est mesuré. Le
tableau~\ref{tab:budget_cifar} donne le budget réel au point $1/12$ ($k = \num{256}$ symboles).

\begin{table}[H]
\centering
\caption{Budget utile par MCS pour une image CIFAR-10 à $k/n = 1/12$}
\label{tab:budget_cifar}
\small
\begin{tabular}{@{}lrrr@{}}
\toprule
\textbf{MCS} & \textbf{bits/symbole} & \textbf{bits utiles} & \textbf{octets} \\
\midrule
QPSK 1/2   & 1,0 & 248  & 31  \\
QPSK 3/4   & 1,5 & 376  & 47  \\
16-QAM 1/2 & 2,0 & 504  & 63  \\
16-QAM 3/4 & 3,0 & 760  & 95  \\
64-QAM 2/3 & 4,0 & 1016 & 127 \\
64-QAM 3/4 & 4,5 & 1136 & 142 \\
\bottomrule
\end{tabular}
\end{table}

Or le plus petit JPEG possible d'une image $32\times32$ pèse de \num{301} à \num{317}~octets : deux à dix fois
le budget. La cause est structurelle — marqueurs et tables de Huffman coûtent environ \num{300}~octets avant
le premier pixel, négligeables en $768\times512$ mais dominants en $32\times32$.

La chaîne emploie donc un auto-encodeur dont le latent est quantifié sur 4~bits (\srcfile{DigitalAE}),
entraîné sans canal : chaque canal latent coûte $64\times4 = \num{256}$~bits, soit 32~octets, et le budget de
142~octets en autorise 4. Ce choix \emph{avantage} délibérément la chaîne séparée.

\section{Codage de canal et réception}
\label{sec:canal_num}

Le train binaire est découpé en paquets d'au plus \num{1024}~bits utiles, terminés par 6~bits de queue. Chaque
paquet traverse : code convolutif $K=7$ poinçonné $\rightarrow$ QAM de Gray $\rightarrow$ canal $\rightarrow$
LLR max-log $\rightarrow$ Viterbi $\rightarrow$ décodage de source.

Rien n'est approché par une formule de taux d'erreur. Un fichier JPEG refusé par le décodeur devient une image
grise uniforme : le récepteur n'a rien d'autre à afficher.

\section{Trois formes de référence}
\label{sec:bornes}

\begin{description}[leftmargin=1.4cm, style=nextline]
  \item[MCS fixe] couple choisi d'avance : plafonne au-dessus du seuil, s'effondre en dessous.
  \item[Meilleur MCS] émetteur « génie » qui connaîtrait le SNR : meilleure adaptation de lien réalisable.
  \item[Borne de capacité] $k\log_2(1+\mathrm{SNR})$ bits \emph{sans erreur} : borne supérieure, pas un
        système constructible.
\end{description}

\begin{keybox}[Pourquoi encadrer la référence]
Battre une référence mal implémentée ne prouverait rien. En publiant les trois formes, on sait exactement ce
qui est dépassé : le chapitre~\ref{chap:resultats} montre que les systèmes appris dépassent largement la
chaîne réelle, dépassent la borne de capacité en précision sur CIFAR-10, mais restent en deçà de cette borne
en PSNR sur grandes images.
\end{keybox}
